%!TeX root=../gardentop.tex
\chapter{»I Shall Live Forever—And Ever—And Ever!«}
	
\begin{figure}[t!]
\centering
\includegraphics[width=\linewidth]{january}
\end{figure}

	\lettrine[lines=6]{B}{ut} they were obliged to wait more than a week because first there came some very windy days and then Colin was threatened with a cold, which two things happening one after the other would no doubt have thrown him into a rage but that there was so much careful and mysterious planning to do and almost every day Dickon came in, if only for a few minutes, to talk about what was happening on the moor and in the lanes and hedges and on the borders of streams. The things he had to tell about otters' and badgers' and water-rats' houses, not to mention birds' nests and field-mice and their burrows, were enough to make you almost tremble with excitement when you heard all the intimate details from an animal charmer and realized with what thrilling eagerness and anxiety the whole busy underworld was working.

»They're same as us,« said Dickon, »only they have to build their homes every year. An' it keeps 'em so busy they fair scuffle to get 'em done.«

The most absorbing thing, however, was the preparations to be made before Colin could be transported with sufficient secrecy to the garden. No one must see the chair-carriage and Dickon and Mary after they turned a certain corner of the shrubbery and entered upon the walk outside the ivied walls. As each day passed, Colin had become more and more fixed in his feeling that the mystery surrounding the garden was one of its greatest charms. Nothing must spoil that. No one must ever suspect that they had a secret. People must think that he was simply going out with Mary and Dickon because he liked them and did not object to their looking at him. They had long and quite delightful talks about their route. They would go up this path and down that one and cross the other and go round among the fountain flower-beds as if they were looking at the »bedding-out plants« the head gardener, Mr~Roach, had been having arranged. That would seem such a rational thing to do that no one would think it at all mysterious. They would turn into the shrubbery walks and lose themselves until they came to the long walls. It was almost as serious and elaborately thought out as the plans of march made by great generals in time of war.

Rumors of the new and curious things which were occurring in the invalid's apartments had of course filtered through the servants' hall into the stable yards and out among the gardeners, but notwithstanding this, Mr~Roach was startled one day when he received orders from Master Colin's room to the effect that he must report himself in the apartment no outsider had ever seen, as the invalid himself desired to speak to him.

»Well, well,« he said to himself as he hurriedly changed his coat, »what's to do now? His Royal Highness that wasn't to be looked at calling up a man he's never set eyes on.«

Mr~Roach was not without curiosity. He had never caught even a glimpse of the boy and had heard a dozen exaggerated stories about his uncanny looks and ways and his insane tempers. The thing he had heard oftenest was that he might die at any moment and there had been numerous fanciful descriptions of a humped back and helpless limbs, given by people who had never seen him.

»Things are changing in this house, Mr~Roach,« said Mrs~Medlock, as she led him up the back staircase to the corridor on to which opened the hitherto mysterious chamber.

»Let's hope they're changing for the better, Mrs~Medlock,« he answered.

»They couldn't well change for the worse,« she continued; »and queer as it all is there's them as finds their duties made a lot easier to stand up under. Don't you be surprised, Mr~Roach, if you find yourself in the middle of a menagerie and Martha Sowerby's Dickon more at home than you or me could ever be.«

There really was a sort of Magic about Dickon, as Mary always privately believed. When Mr~Roach heard his name he smiled quite leniently.

»He'd be at home in Buckingham Palace or at the bottom of a coal mine,« he said. »And yet it's not impudence, either. He's just fine, is that lad.«

It was perhaps well he had been prepared or he might have been startled. When the bedroom door was opened a large crow, which seemed quite at home perched on the high back of a carven chair, announced the entrance of a visitor by saying »Caw—Caw« quite loudly. In spite of Mrs~Medlock's warning, Mr~Roach only just escaped being sufficiently undignified to jump backward.

The young Rajah was neither in bed nor on his sofa. He was sitting in an armchair and a young lamb was standing by him shaking its tail in feeding-lamb fashion as Dickon knelt giving it milk from its bottle. A squirrel was perched on Dickon's bent back attentively nibbling a nut. The little girl from India was sitting on a big footstool looking on.

»Here is Mr~Roach, Master Colin,« said Mrs~Medlock.

The young Rajah turned and looked his servitor over—at least that was what the head gardener felt happened.

»Oh, you are Roach, are you?« he said. »I sent for you to give you some very important orders.«

»Very good, sir,« answered Roach, wondering if he was to receive instructions to fell all the oaks in the park or to transform the orchards into water-gardens.

»I am going out in my chair this afternoon,« said Colin. »If the fresh air agrees with me I may go out every day. When I go, none of the gardeners are to be anywhere near the Long Walk by the garden walls. No one is to be there. I shall go out about two o'clock and every one must keep away until I send word that they may go back to their work.«

»Very good, sir,« replied Mr~Roach, much relieved to hear that the oaks might remain and that the orchards were safe.

»Mary,« said Colin, turning to her, »what is that thing you say in India when you have finished talking and want people to go?«

»You say, 'You have my permission to go,'« answered Mary.

The Rajah waved his hand.

»You have my permission to go, Roach,« he said. »But, remember, this is very important.«

»Caw—Caw!« remarked the crow hoarsely but not impolitely.

»Very good, sir. Thank you, sir,« said Mr~Roach, and Mrs~Medlock took him out of the room.

Outside in the corridor, being a rather good-natured man, he smiled until he almost laughed.

»My word!« he said, »he's got a fine lordly way with him, hasn't he? You'd think he was a whole Royal Family rolled into one—Prince Consort and all.«

»Eh!« protested Mrs~Medlock, »we've had to let him trample all over every one of us ever since he had feet and he thinks that's what folks was born for.«

»Perhaps he'll grow out of it, if he lives,« suggested Mr~Roach.

»Well, there's one thing pretty sure,« said Mrs~Medlock. »If he does live and that Indian child stays here I'll warrant she teaches him that the whole orange does not belong to him, as Susan Sowerby says. And he'll be likely to find out the size of his own quarter.«

Inside the room Colin was leaning back on his cushions.

»It's all safe now,« he said. »And this afternoon I shall see it—this afternoon I shall be in it!«

Dickon went back to the garden with his creatures and Mary stayed with Colin. She did not think he looked tired but he was very quiet before their lunch came and he was quiet while they were eating it. She wondered why and asked him about it.

»What big eyes you've got, Colin,« she said. »When you are thinking they get as big as saucers. What are you thinking about now?«

»I can't help thinking about what it will look like,« he answered.

»The garden?« asked Mary.

»The springtime,« he said. »I was thinking that I've really never seen it before. I scarcely ever went out and when I did go I never looked at it. I didn't even think about it.«

»I never saw it in India because there wasn't any,« said Mary.

Shut in and morbid as his life had been, Colin had more imagination than she had and at least he had spent a good deal of time looking at wonderful books and pictures.

»That morning when you ran in and said 'It's come! It's come!' you made me feel quite queer. It sounded as if things were coming with a great procession and big bursts and wafts of music. I've a picture like it in one of my books—crowds of lovely people and children with garlands and branches with blossoms on them, every one laughing and dancing and crowding and playing on pipes. That was why I said, 'Perhaps we shall hear golden trumpets' and told you to throw open the window.«

»How funny!« said Mary. »That's really just what it feels like. And if all the flowers and leaves and green things and birds and wild creatures danced past at once, what a crowd it would be! I'm sure they'd dance and sing and flute and that would be the wafts of music.«

They both laughed but it was not because the idea was laughable but because they both so liked it.

A little later the nurse made Colin ready. She noticed that instead of lying like a log while his clothes were put on he sat up and made some efforts to help himself, and he talked and laughed with Mary all the time.

»This is one of his good days, sir,« she said to Dr~Craven, who dropped in to inspect him. »He's in such good spirits that it makes him stronger.«

»I'll call in again later in the afternoon, after he has come in,« said Dr~Craven. »I must see how the going out agrees with him. I wish,« in a very low voice, »that he would let you go with him.«

»I'd rather give up the case this moment, sir, than even stay here while it's suggested,« answered the nurse with sudden firmness.

»I hadn't really decided to suggest it,« said the doctor, with his slight nervousness. »We'll try the experiment. Dickon's a lad I'd trust with a new-born child.«

The strongest footman in the house carried Colin down-stairs and put him in his wheeled chair near which Dickon waited outside. After the manservant had arranged his rugs and cushions the Rajah waved his hand to him and to the nurse.

»You have my permission to go,« he said, and they both disappeared quickly and it must be confessed giggled when they were safely inside the house.

Dickon began to push the wheeled chair slowly and steadily. Mistress Mary walked beside it and Colin leaned back and lifted his face to the sky. The arch of it looked very high and the small snowy clouds seemed like white birds floating on outspread wings below its crystal blueness. The wind swept in soft big breaths down from the moor and was strange with a wild clear scented sweetness. Colin kept lifting his thin chest to draw it in, and his big eyes looked as if it were they which were listening—listening, instead of his ears.

»There are so many sounds of singing and humming and calling out,« he said. »What is that scent the puffs of wind bring?«

»It's gorse on th' moor that's openin' out,« answered Dickon. »Eh! th' bees are at it wonderful to-day.«

Not a human creature was to be caught sight of in the paths they took. In fact every gardener or gardener's lad had been witched away. But they wound in and out among the shrubbery and out and round the fountain beds, following their carefully planned route for the mere mysterious pleasure of it. But when at last they turned into the Long Walk by the ivied walls the excited sense of an approaching thrill made them, for some curious reason they could not have explained, begin to speak in whispers.

»This is it,« breathed Mary. »This is where I used to walk up and down and wonder and wonder.«

»Is it?« cried Colin, and his eyes began to search the ivy with eager curiousness. »But I can see nothing,« he whispered. »There is no door.«

»That's what I thought,« said Mary.

Then there was a lovely breathless silence and the chair wheeled on.

»That is the garden where Ben Weatherstaff works,« said Mary.

»Is it?« said Colin.

A few yards more and Mary whispered again.

»This is where the robin flew over the wall,« she said.

»Is it?« cried Colin. »Oh! I wish he'd come again!«

»And that,« said Mary with solemn delight, pointing under a big lilac bush, »is where he perched on the little heap of earth and showed me the key.«

Then Colin sat up.

»Where? Where? There?« he cried, and his eyes were as big as the wolf's in Red Riding-Hood, when Red Riding-Hood felt called upon to remark on them. Dickon stood still and the wheeled chair stopped.

»And this,« said Mary, stepping on to the bed close to the ivy, »is where I went to talk to him when he chirped at me from the top of the wall. And this is the ivy the wind blew back,« and she took hold of the hanging green curtain.

»Oh! is it—is it!« gasped Colin.

»And here is the handle, and here is the door. Dickon push him in—push him in quickly!«

And Dickon did it with one strong, steady, splendid push.

But Colin had actually dropped back against his cushions, even though he gasped with delight, and he had covered his eyes with his hands and held them there shutting out everything until they were inside and the chair stopped as if by magic and the door was closed. Not till then did he take them away and look round and round and round as Dickon and Mary had done. And over walls and earth and trees and swinging sprays and tendrils the fair green veil of tender little leaves had crept, and in the grass under the trees and the gray urns in the alcoves and here and there everywhere were touches or splashes of gold and purple and white and the trees were showing pink and snow above his head and there were fluttering of wings and faint sweet pipes and humming and scents and scents. And the sun fell warm upon his face like a hand with a lovely touch. And in wonder Mary and Dickon stood and stared at him. He looked so strange and different because a pink glow of colour had actually crept all over him—ivory face and neck and hands and all.

»I shall get well! I shall get well!« he cried out. »Mary! Dickon! I shall get well! And I shall live forever and ever and ever!«